% Shared visual language for the seven analytical papers.
% This file defines marks, not layouts: every figure must still choose a grammar
% from its reader question.  The duplicate under whitepaper/figures is kept in
% lockstep for standalone Chapter I/II builds.
\usetikzlibrary{matrix,patterns.meta,shapes.geometric}

% ===========================================================================
% THE TYPOGRAPHIC LAW
% ===========================================================================
% Companion to the line-weight law in figures/pd-palette.tex (1px texture,
% 1.5px linework, 2px enclosure). Same shape of rule: name the roles, fix the
% values, one place.
%
% ONE SIZE. Every named text role in a figure is \pdfiglabelsize
% (= \footnotesize). Roles separate by WEIGHT, SLOPE, FAMILY and INK -- never
% by size. A fragment never writes a size command of its own; if a distinction
% needs making, it takes the role that makes it.
%
% Where the size comes from. Measured, not chosen:
%   \tiny  \scriptsize  \footnotesize  \small  \normalsize
%   5.23    6.97          8.72           9.55    10.46  pt   Book (10.5pt
%                                                            Pagella, 4.5in
%                                                            column, 7x10in)
%   5.98    7.97          8.97           9.96    10.91  pt   standalone chapter
%                                                            (11pt Latin Modern)
% \scriptsize renders at 6.97pt in the Book -- under figcheck's 7pt T1 floor
% before a single width promotion, so it is illegal by measurement, not taste.
% \footnotesize holds 8.72pt and still clears the floor at the 0.85 \resizebox
% limit craft-rules sets (7.41pt). It is also what the corpus already voted
% for: 78 of the 91 local font= declarations this law replaces were already
% \footnotesize, \footnotesize\bfseries or \footnotesize\itshape.
%
% \pdfigbasesize (= \small) is the tikzpicture's ambient size only, so an
% unstyled node is never SMALLER than a styled one; every named role steps
% down exactly one notch from it and none steps down twice.
%
% ONE FACE, AND IT IS THE DOCUMENT'S. \pdfiglabelfamily is empty by default,
% which means figure type inherits whatever face the document is setting --
% Palatino in the Book, Computer Modern in a standalone chapter. Figures
% agreeing with EACH OTHER is not the goal; figures agreeing with the
% paragraph beside them is. A fragment therefore never names a family. Not
% \sffamily (which resolves to Latin Modern Sans against a Computer Modern
% page and reads as a foreign object on it), not \rmfamily, not
% \fontfamily{...}. P18 in tikz_precheck.py is an error for exactly that.
% An edition MAY substitute a face -- the Swiss and technical editions set
% \pdfiglabelfamily to \pdgrotesk, because grotesk display is their declared
% character -- and it does so in ONE command in ONE file, so the substitution
% is a decision on the record rather than an accident in a fragment. The
% only family a fragment may ask for is the identifier one, and it asks for
% it by role (pd mono label / \pdfiglabelmono), never by \ttfamily.
%
% What the roles are, and what each one is for:
%   pd panel title        bold        the name of a panel
%   pd row label          bold        an actor/row name east of its lane
%   pd reverse row label  bold, cream a row name sitting on a dark band
%   pd axis label         regular     an axis title or a tick numeral
%   pd direct label       regular     a value or name riding its own mark
%   pd reverse label      regular, cream  a direct label on a dark band
%   pd mono label         mono        an identifier, path, or literal
%   pd verdict            bold        a terminal outcome stamp (CLEAR, REJECT)
%   pd note               italic      one aside per figure
%   pd legend             regular     a pgfplots legend entry
%   pd state / terminal / actor / artifact / decision   node text
%   pd axis               the whole pgfplots typographic block, one key
%   \pdfigsub             a gloss line under a heading INSIDE a node
%   \pdfigmath            math carrying a sub/superscript, inside a label
%
% A fragment that needs a compound style of its own composes it from a role
% above (`gate/.style={pd decision,minimum width=2cm}`), or, where no role
% fits, from the \pdfiglabel* handles here. It never writes a raw size
% command: P16 in skills/harbor-chartwork/scripts/tikz_precheck.py is an
% error for exactly that, and P10/P11 have been errors for \tiny and
% \scriptsize since before this law.
%
% What does NOT transfer from the screen rules these roles otherwise follow.
% The dataviz method those rules come from assumes an interactive chart: a
% tooltip carries any value a label could not fit, a filter row rescopes the
% view, hover lifts the hovered mark. A printed page has none of that, and the
% absence is not a defect to design around. The INTENT of the tooltip rule --
% "every value is reachable without gating" -- does transfer, and on this page
% it is discharged somewhere else: the value is in the caption, in the margin
% apparatus, or in the chapter's own worked example. So a label that will not
% fit inside its mark moves outside the mark or is dropped to the caption; it
% is never shrunk to fit, because there is no size below \footnotesize to shrink
% to, and never clipped. Do not add a hover affordance, a legend toggle, or a
% filter to a figure in this book.
%
% IDENTIFIERS GET ONE SPELLING. A key, a cell id, a file name or any other
% literal a reader could type is set in `pd mono label` (or \pdfiglabelmono
% inside a compound style), spelled exactly as the code spells it: sk_A, not
% skA and not sk with a subscript A; card0, not card-nought and not a Unicode
% subscript zero. A figure that shows the same identifier two ways has told
% the reader they are two things. The subscripted form belongs to MATHEMATICS
% ($\kappa_7$, an index into a sequence) and the literal form to CODE -- one
% figure may use both, for different objects, but never both for the same
% one. (A precheck for two spellings of one identifier inside a fragment is
% being added on another branch; this is the convention it will enforce.)
%
% SHAPE IS A CHANNEL, AND IT HAS TO SURVIVE THE PRINT SIZE. pd datum, pd
% focus datum and pd caution datum carry three different forms, because the
% house palette is three low-chroma print inks and every distinction it makes
% has to be doubled by something that is not colour. That only counts if the
% form is still legible at 2.1-2.7 pt: an edition that flattens all three to
% one form has removed the second channel and left the figure separating its
% classes by hue alone, which fails in greyscale and for a colour-blind
% reader. Any edition override that restyles a datum must keep three
% distinguishable forms, and must be checked on a render at 1.0x, not
% reasoned about from the point sizes.
%
% One screen rule that does NOT bind here, and why. "Text never wears the data
% colour" exists because a light categorical hue is illegible as text on a
% light surface. The house figure hues are print inks, not screen hues --
% hhink #1B1712, hhteal #00564C, hhamber #6B4500 -- and all three clear text
% contrast on cream, so a label MAY wear its series' ink to bind itself to a
% curve, which is direct labelling and saves a legend. The rule that survives
% is the reason behind it: text wears one of those three inks at full
% strength, or hhink!80!hhgray for recessive furniture, and nothing else. A
% tint of an ink (text=hhgray, text=hhteal!60) is not a role; it is a fade.
% ===========================================================================

\providecommand{\pdfiglabelfamily}{}            % an edition may set \pdgrotesk
\providecommand{\pdfiglabelsize}{\footnotesize} % every named text role
\providecommand{\pdfigbasesize}{\small}         % the picture's ambient size
\providecommand{\pdfiglabel}{\pdfiglabelfamily\pdfiglabelsize}
\providecommand{\pdfiglabelbold}{\pdfiglabel\bfseries}
\providecommand{\pdfiglabelitalic}{\pdfiglabel\itshape}
\providecommand{\pdfiglabelmono}{\pdfiglabelsize\ttfamily}
% The one role that is not a node style, because it lives INSIDE a node's text:
% the second line under a heading that glosses it ("reviewer role / obligation +
% capability + authority"). Every fragment that needed it reached for \tiny or
% \scriptsize, which is how seven \tiny and twelve \scriptsize survived P10/P11
% as hard errors in the corpus. It steps down by SLOPE and WEIGHT, not size, so
% it stays legible and stays legal, and it touches neither family nor ink --
% which is what lets it sit on a reversed node without vanishing.
\providecommand{\pdfigsub}{\mdseries\itshape}
% The one place the law's own size is not the right size, with the measurement
% that says so. A subscript is set at roughly 70% of its base, so math with a
% subscript inside a \pdfiglabelsize label renders its subscript at 5.98pt in a
% standalone chapter and 6.36pt in the Book -- under figcheck's 7pt T1 floor,
% measured, not estimated. One notch up puts the same subscript at 7.97pt and
% 7.64pt, over the floor, while the base glyph stays close enough to its
% neighbours that the label still reads as one label. Use it ONLY for a math
% group that carries a sub- or superscript; plain math takes the label's own
% size like any other text.
\providecommand{\pdfigmath}{\normalsize}

% ---------------------------------------------------------------------------
% THE GUIDE DASH, IN ABSOLUTE POINTS, IN ONE PLACE.
%
% `pd guide` used to read `line width=.45pt,densely dotted`. `densely dotted`
% expands to `dash pattern=on \pgflinewidth off 1pt`, and \pgflinewidth is read
% when the KEY is processed, not when the path is stroked. So the Swiss
% edition's `pd guide/.append style={draw=hhink!40,line width=.5pt}` widened the
% stroke and left the dash at the .45pt already baked into it. The canonical
% Book therefore shipped a 0.448pt dot on a 0.498pt stroke.
%
% Measured on the Book (Swiss, the built edition) at 1.0x / 150 dpi:
%   on 0.448pt = 0.93 device px -- UNDER ONE PIXEL
%   off 0.996pt = 2.08 px, period 3.01 px
%   peak ink per dot swings 0.173-0.345 (a 2.0x range) purely with where the
%     dot lands on the pixel grid; the line's weight is a rasteriser accident
%   the darkest pixel of the guide (0.30) is DARKER than every pixel of the
%     solid `pd hairline` (0.25/0.19) beside it, while its mean ink (2.3%) is a
%     quarter of the hairline's (8.9%) -- the two lines are not separated on any
%     axis a reader can name
%   declared separation from `pd hairline`: hhink!40 against hhink!45, 4.6% of
%     the grey scale
%
% So a fragment that used `pd guide` to mean "this line is not that line" was
% relying on ink that does not arrive on the page. Twenty-five of the corpus's
% thirty-six guide draws carry exactly that meaning.
%
% The replacement is stated absolutely and measures phase-STABLE: every dot
% renders at the same weight wherever it falls on the grid (peak column ink
% 0.808, invariant across twelve sub-pixel offsets). It is a macro rather than
% three numbers because an edition that restates one of them and not the others
% is how this defect happened. The geometry is a LEGIBILITY FLOOR, not an
% edition's character: an edition may move the guide's ink, and states the
% geometry through these two macros so it cannot half-restate it.
%
% Any dash in this corpus is stated in absolute points. The `dotted` family
% (`dotted`, `densely dotted`, `loosely dotted`) derives its on-length from the
% line width and is therefore banned -- P23 in tikz_precheck.py. `dashed` and
% its relatives are absolute (on 3pt off 3pt) and are fine.
\providecommand{\pdguidewidth}{.7pt}
% Three lengths, not one macro holding the whole key: pgfkeys splits a key list
% on `,` and `=` before expanding anything (so a macro carrying `dash
% pattern=...` is read as one key NAME), and TikZ's dash-pattern scanner reads
% the literal words `on`/`off` without expanding a macro that hides them. A
% dimension after each word does expand, which is the seam that works.
\providecommand{\pdguideon}{1.2pt}
\providecommand{\pdguideoff}{2.0pt}

\tikzset{
  pd figure/.style={font=\pdfiglabelfamily\pdfigbasesize,text=hhink},
  pd hairline/.style={draw=hhgray!78,line width=.5pt},
  pd rule/.style={draw=hhink,line width=.62pt},
  pd focus rule/.style={draw=hhteal,line width=1.05pt},
  pd caution rule/.style={draw=hhamber,line width=1.05pt},
  % Darker per dot than the hairline, lighter in mean ink than it: a guide is a
  % construction line the reader traces from a mark to an axis, so each dot has
  % to be a definite mark, while the line as a whole has to recede behind the
  % data. The hairline is the opposite -- continuous and faint. Those are the
  % two axes that now separate them, instead of 4.6% of the grey scale.
  pd guide/.style={draw=hhink!70,line width=\pdguidewidth,dash pattern=on \pdguideon off \pdguideoff},
  % A LATTICE is the rule work a figure stands on: column separators, row
  % baselines, a time grid. It is NOT dotted, and that is the point. Dotted
  % versus solid is a semantic channel -- it distinguishes a construction line
  % from a drawn edge -- and thirteen fragments were spending it on background
  % grid where it distinguishes nothing. A dash that carries no meaning is the
  % first ink a rasteriser is entitled to lose, and it did. A lattice is
  % continuous and quiet; a guide is discontinuous and definite. Use `pd guide`
  % only where the reader must see that the line is not a drawn edge.
  pd lattice/.style={draw=hhgray!52,line width=.4pt},
  pd arrow/.style={-{Stealth[length=2.0mm,width=1.25mm]},draw=hhink,line width=.68pt},
  pd focus arrow/.style={-{Stealth[length=2.0mm,width=1.25mm]},draw=hhteal,line width=.88pt},
  pd caution arrow/.style={-{Stealth[length=2.0mm,width=1.25mm]},draw=hhamber,line width=.88pt,dashed},
  pd row label/.style={font=\pdfiglabelbold,text=hhink,anchor=east},
  pd panel title/.style={font=\pdfiglabelbold,text=hhink,align=left},
  pd axis label/.style={font=\pdfiglabel,text=hhink!80!hhgray},
  % inner ysep is deliberately larger than inner xsep: at 1.5pt the descenders
  % of a last line (q, y, p, g) reach the bottom of the cream backing under
  % Computer Modern, where they clear it under Palatino. A label should not
  % need its own patch to hold its own tails, and the sensitivity is the
  % face's, not the fragment's.
  pd direct label/.style={font=\pdfiglabel,text=hhink,align=left,
    inner xsep=1.5pt,inner ysep=2.5pt,fill=hhpaper},
  pd note/.style={font=\pdfiglabelitalic,text=hhink!80!hhgray,align=left},
  pd tick/.style={draw=hhgray,line width=.45pt},
  pd datum/.style={circle,fill=hhink,draw=none,inner sep=2.1pt},
  pd focus datum/.style={circle,fill=hhteal,draw=hhpaper,line width=.55pt,inner sep=2.5pt},
  pd caution datum/.style={diamond,fill=hhamber,draw=hhpaper,line width=.55pt,inner sep=2.5pt},
  pd state/.style={draw=hhink,line width=.62pt,fill=hhsand!30,rounded corners=1.5pt,
    font=\pdfiglabel,inner xsep=8pt,inner ysep=6pt,align=center,minimum height=9mm},
  pd terminal/.style={pd state,double,double distance=1.2pt},
  pd actor/.style={draw=hhink,line width=.62pt,fill=hhpaper,rounded corners=1.5pt,
    font=\pdfiglabel,inner xsep=8pt,inner ysep=5pt,align=center},
  pd artifact/.style={draw=hhgray!80,line width=.55pt,fill=hhsand!30,
    font=\pdfiglabel,inner xsep=8pt,inner ysep=5pt,align=center},
  pd boundary/.style={draw=hhgray!70,line width=.55pt,dashed,rounded corners=2pt},
  % These three are pale washes: they are frequently layered on top of `pd
  % state` (directly, or via a compound style like `cell`), whose swiss-edition
  % override (figures/pd-figure-language-swiss.tex) appends text=hhpaper for the
  % common case of a solid, saturated fill. A pale wash is the opposite case --
  % text must stay dark ink regardless of what came before it in the style
  % list, so each one restates text=hhink explicitly rather than relying on
  % whatever was inherited.
  pd focus fill/.style={fill=hhteal!24,draw=hhteal!65,line width=.4pt,text=hhink},
  pd caution fill/.style={fill=hhamber!26,draw=hhamber!70,line width=.4pt,text=hhink},
  pd neutral fill/.style={fill=hhsand!40,draw=hhgray!65,line width=.4pt,text=hhink},
  pd ink fill/.style={fill=hhink,draw=hhink,line width=.4pt,text=hhpaper},
  pd hatch/.style={pattern={Lines[angle=45,distance=4pt,line width=.28pt]},pattern color=hhgray!55},
  % --- roles the corpus was inventing locally, now named -------------------
  % Each is built ON TOP of the role it varies, so an edition's .append style
  % on the parent reaches it without the edition file having to list it.
  pd decision/.style={pd state,diamond,aspect=1.55,font=\pdfiglabelbold,
    minimum height=0pt,inner xsep=3pt,inner ysep=3pt},
  % A shade more side bearing than its parent: a monospace glyph box overhangs
  % its own advance width, so the cream backing at the parent's 1.5pt clipped
  % the last character of a long identifier.
  pd mono label/.style={pd direct label,font=\pdfiglabelmono,inner xsep=2.5pt,inner ysep=2.5pt},
  pd verdict/.style={pd direct label,font=\pdfiglabelbold},
  pd reverse label/.style={pd direct label,text=hhpaper,fill=none},
  pd reverse row label/.style={pd row label,text=hhpaper},
  pd legend/.style={font=\pdfiglabel,text=hhink,draw=none,fill=none},
  % A KIND TAG names what a row or column IS, rather than what it says: an
  % audience ("for the operator"), a category, a running head. craft-rules S6
  % already reserves small caps for exactly this, and small caps is the
  % smallest effective difference that marks a change in kind -- it needs no
  % size drop, no second colour and no box.
  pd kind tag/.style={pd axis label,font=\pdfiglabel\scshape,align=right},
  % A BADGE is a short boxed fact riding beside a row -- a part and chapter
  % range, a version, a count. Hairline edge, no fill: it is a label with a
  % boundary, not a region.
  pd badge/.style={pd hairline,font=\pdfiglabel,text=hhink,align=center,
    inner sep=2.4pt,fill=none},
}

% The pgfplots half of the same law. Before this existed, every plot fragment
% restated title style/label style/tick label style by hand, and they disagreed
% -- four said the tick numerals were hhink!80!hhgray, three said hhgray, and
% the axis rule was drawn four different ways. `pd axis` is the one place now.
% The axis line is `pd hairline` because an axis is reference furniture, not
% data: heavy for data, light for the frame that carries it.
\ifdefined\pgfplotsset
\pgfplotsset{
  pd axis/.style={
    title style={font=\pdfiglabelbold,text=hhink,yshift=-2pt},
    label style={font=\pdfiglabel,text=hhink!80!hhgray},
    tick label style={font=\pdfiglabel,text=hhink!80!hhgray},
    legend style={pd legend},
    axis line style={pd hairline},
    tick style={pd tick},
  },
}
\fi

% Small reusable marks.  They deliberately avoid prose-bearing boxes.
\newcommand{\pdfigaxis}[4]{%
  \draw[pd rule,-{Stealth[length=1.65mm,width=1.05mm]}] (#1,#2) -- (#3,#2);
  \node[pd axis label,anchor=west] at (#3,#2) {#4};
}
\newcommand{\pdfigvaxis}[4]{%
  \draw[pd rule,-{Stealth[length=1.65mm,width=1.05mm]}] (#1,#2) -- (#1,#3);
  \node[pd axis label,anchor=south,align=center] at (#1,#3) {#4};
}

% ---------------------------------------------------------------------------
% Editions. The Book sets \pdedition (maritime | swiss | technical) and the
% matching override file restyles the same names, so one drawing renders in
% each edition's character. Standalone chapters have no edition. An edition
% changes the FAMILY (via \pdfiglabelfamily) and the character of the marks;
% it never changes the size, because the size is the law above.
% ---------------------------------------------------------------------------
\ifdefined\pdedition\IfFileExists{figures/pd-figure-language-\pdedition.tex}{\input{figures/pd-figure-language-\pdedition}}{}\fi
